\section{Evaluation}
\label{sec:eval}

All measurements are on a single Grace-class AArch64 host (Cortex-X925,
$20$ cores, $121$\,GB RAM, Linux). Each server is driven over its native
protocol against a $10$M-row \texttt{MergeTree} table
(\texttt{id UInt64, user\_id UInt32, ts DateTime, value Float64, category
LowCardinality(String)}). We report steady-state resident set (RSS, from
\texttt{/proc/<pid>/status}) and queries-per-second. ``Official'' is the
upstream binary, which ships PGO+LTO+BOLT.

\subsection{The fork costs nothing}

Built with identical flags ($-O3$, no PGO/LTO), Datastore and upstream are
statistically indistinguishable: RSS within $1$\,MB ($143$ vs $144$\,MB) and
throughput within noise across aggregation, top-$k$, and filter at concurrencies
$1$ through $64$. This establishes that every subsequent gain is attributable to
the optimization pipeline, not to the fork.

\subsection{Effect of each optimization stage}

Table~\ref{tab:stages} reports resident set and binary size as the pipeline is
applied. ThinLTO alone already brings the resident set below the official
binary's; PGO halves it again via hot/cold splitting; removing XRay and applying
BOLT delivers the final size and a further locality gain.

\begin{table}[H]
\centering
\small
\begin{tabular}{@{}lrr@{}}
\toprule
Configuration & RSS & Binary \\
\midrule
Datastore $-O3$            & $143$\,MB & $709$\,MB \\
\quad + ThinLTO            & $113$\,MB & $684$\,MB \\
\quad + PGO                & $53$\,MB  & --- \\
\quad + no-XRay, + BOLT    & $\mathbf{49}$--$\mathbf{50}$\,MB & $\mathbf{648}$\,MB \\
\midrule
Official (PGO+LTO+BOLT)    & $116$\,MB & --- \\
\bottomrule
\end{tabular}
\caption{Resident set and binary size by optimization stage. ThinLTO alone
($113$\,MB) already beats the official binary's $116$\,MB; the full pipeline
reaches $49$--$50$\,MB, $\sim\!57\%$ below official.}
\label{tab:stages}
\end{table}

\subsection{Throughput at parity memory}

At equal correctness the optimized Datastore matches or beats upstream on the
read mix. On the stable mid-concurrency point ($c{=}4$, $10$M rows) the
no-XRay PGO+LTO+BOLT build leads upstream on filters by $25$--$40\%$
($397$ vs $284$ QPS), on top-$k$ by $10$--$30\%$, and tracks or exceeds it on
aggregation, while using less than half the resident memory. We caution that
single-shot QPS on a shared host carries contention noise; the memory result,
by contrast, is invariant to load and reproduces on every run.

\subsection{Memory is flat in connection count}

To isolate per-connection cost we issue a trivial query across $10$ to
$10{,}000$ concurrent clients (Table~\ref{tab:scale}). Resident set is flat:
$10{,}000$ connections cost the same memory as $10$. This is the bounded-state
connection model --- the server pools threads and holds negligible per-connection
state --- and Datastore is $\sim\!15$\,MB leaner than upstream at every point.

\begin{table}[H]
\centering
\small
\begin{tabular}{@{}lrr@{}}
\toprule
Concurrent clients & Datastore RSS & Official RSS \\
\midrule
$10$      & $114$\,MB & $129$\,MB \\
$100$     & $113$\,MB & $129$\,MB \\
$1{,}000$ & $113$\,MB & $128$\,MB \\
$10{,}000$& $\mathbf{113}$\,MB & $128$\,MB \\
\bottomrule
\end{tabular}
\caption{Resident set vs.\ concurrent connection count. Per-connection memory is
effectively zero; $10{,}000$ clients cost the same as $10$. (Measured on an
LTO build before the PGO/BOLT stages of Table~\ref{tab:stages}.)}
\label{tab:scale}
\end{table}

\subsection{ZAP pipelining flattens dependent-chain latency}

We measure end-to-end latency of a dependent chain of $n$ calls --- each input
the previous output --- at a $0.5$\,ms round trip, comparing ZAP promise
pipelining against a request-response baseline. The baseline grows linearly in
$n$; ZAP is flat, since the whole chain resolves in one logical round trip. At
depth $64$ this is a $63\times$ reduction in latency. This is the property that
motivates ZAP for analytics clients that issue dependent queries (build a
filtered view, then aggregate it, then top-$k$ it) rather than independent ones.
